\chapter{Nhanh và Chậm}
\markboth{Nhanh và Chậm}{Fast and Slow}

\section{Một người muốn thành công nhanh nên bỏ qua nền móng.}
\begin{truyen}{Một người muốn thành công nhanh nên bỏ qua nền móng.}{Speed without foundation collapses early.}
\chuhoa{T}{hấy người khác đi lên nhanh, anh cũng muốn có kết quả thật sớm. Anh lao vào nhiều cơ hội, học lướt, quyết vội, và khoe trước khi hiểu rõ mình đang làm gì. Chỉ vài tháng sau, mọi thứ gãy xuống vì thứ đi nhanh nhất thường là thứ chưa kịp bám rễ.}
\textit{(Seeing others rise quickly, he also wanted fast results. He rushed into opportunities, skimmed instead of learning, decided in haste, and spoke before he understood. A few months later everything collapsed, because what grows fastest is often what has not yet taken root.)}
\end{truyen}
\begin{baihoc}
Nhanh không xấu, nhưng nhanh trên nền yếu rất dễ đổ.
\textit{(Speed is not bad, but speed on weak foundations collapses easily.)}
\end{baihoc}
\begin{ghinhoanh}
\textbf{Short English to remember:}

Speed needs roots.
\end{ghinhoanh}
\begin{tuhocnhanh}
1. Em đang vội ở đâu mà chưa xây chắc nền? \textit{(Where are you rushing without building a strong base?)}

2. Nếu chậm lại một chút, em giữ được điều gì? \textit{(If you slowed down, what could you protect?)}
\end{tuhocnhanh}
\ngancach

\section{Một học sinh tiến bộ chậm nhưng hiểu rất sâu.}
\begin{truyen}{Một học sinh tiến bộ chậm nhưng hiểu rất sâu.}{Slow progress can be deep progress.}
\chuhoa{E}{m không nổi bật ngay từ đầu. Nhưng em kiên trì lấp từng lỗ hổng, làm đi làm lại các bài cơ bản, và hỏi cho ra bản chất. Một năm sau, khi nhiều bạn hụt hơi, em lại đi đều vì cái chậm ban đầu đã đổi thành độ chắc.}
\textit{(The student was not impressive at first. But step by step the gaps were filled, basic exercises were repeated, and true understanding was pursued. A year later, when many others lost momentum, the once-slow learner moved steadily because early slowness had turned into strength.)}
\end{truyen}
\begin{baihoc}
Chậm không phải lúc nào cũng là kém; nhiều khi đó là cách nền móng được đổ sâu.
\textit{(Slow does not always mean weak; often it is how foundations are poured deeper.)}
\end{baihoc}
\begin{ghinhoanh}
\textbf{Short English to remember:}

Slow can become strong.
\end{ghinhoanh}
\begin{tuhocnhanh}
1. Em có đang tự ép mình phải nhanh như người khác không? \textit{(Are you forcing yourself to move as fast as others?)}

2. Nền móng nào của em cần thêm thời gian? \textit{(Which foundation in you still needs more time?)}
\end{tuhocnhanh}
\ngancach

\section{Một mối quan hệ bắt đầu quá nhanh nên cũng rời rất nhanh.}
\begin{truyen}{Một mối quan hệ bắt đầu quá nhanh nên cũng rời rất nhanh.}{Fast closeness is not always deep closeness.}
\chuhoa{C}{ả hai cuốn vào nhau bằng cảm xúc mạnh, nói với nhau những lời rất lớn khi còn chưa hiểu nhau đủ. Đến khi khác biệt hiện ra, mối quan hệ rạn nhanh như lúc nó bắt đầu. Có những điều muốn bền thì phải đi qua một quãng chậm để thật sự hiểu.}
\textit{(They were pulled together by strong feelings and said huge words before they knew each other well. When differences emerged, the relationship cracked as quickly as it had formed. Some things last only if they pass through a slower season of understanding.)}
\end{truyen}
\begin{baihoc}
Nhanh cho cảm xúc không bảo đảm bền cho tình cảm.
\textit{(Fast emotion does not guarantee lasting affection.)}
\end{baihoc}
\begin{ghinhoanh}
\textbf{Short English to remember:}

Fast feelings need slow truth.
\end{ghinhoanh}
\begin{tuhocnhanh}
1. Em có đang muốn điều gì lớn lên quá nhanh không? \textit{(Are you wanting something to grow too quickly?)}

2. Điều gì cần thêm thời gian để thành thật hơn? \textit{(What needs more time to become more genuine?)}
\end{tuhocnhanh}
\ngancach

\section{Một người quyết rất nhanh rồi sửa rất lâu.}
\begin{truyen}{Một người quyết rất nhanh rồi sửa rất lâu.}{Quick choices can demand long repair.}
\chuhoa{A}{nh tự hào vì mình dứt khoát. Nhưng nhiều lần sự dứt khoát ấy chỉ là thiếu kiên nhẫn. Mỗi lần xử lý hậu quả, anh mới thấy không phải cứ ra quyết định nhanh là tiết kiệm thời gian.}
\textit{(He was proud of being decisive. Yet many times that decisiveness was simply impatience. Each time he repaired the consequences, he realized that making quick decisions does not automatically save time.)}
\end{truyen}
\begin{baihoc}
Nhanh lúc chọn mà chậm lúc sửa chưa chắc là khôn.
\textit{(Being fast in choosing but slow in repairing is not necessarily wise.)}
\end{baihoc}
\begin{ghinhoanh}
\textbf{Short English to remember:}

Quick choices can cost time.
\end{ghinhoanh}
\begin{tuhocnhanh}
1. Em có đang quyết quá nhanh ở chỗ cần suy nghĩ thêm không? \textit{(Are you deciding too quickly where more thought is needed?)}

2. Hậu quả nào em đang sửa vì một phút nóng vội? \textit{(What consequence are you repairing because of one rushed moment?)}
\end{tuhocnhanh}
\ngancach

\section{Một người đi chậm nhưng biết dừng để nhìn lại mình.}
\begin{truyen}{Một người đi chậm nhưng biết dừng để nhìn lại mình.}{Slow steps can see more truth.}
\chuhoa{K}{hông phải ai đi chậm cũng tụt lại. Có người chậm vì họ còn đủ tỉnh để hỏi mình đang đi đúng chưa, đang sống tử tế chưa, đang mải chạy vì điều gì. Chính những lần dừng ấy giữ cho con đường không lệch quá xa.}
\textit{(Not everyone who moves slowly is falling behind. Some move slowly because they remain awake enough to ask whether they are on the right road, living decently, and running for the right reasons. Such pauses keep the road from drifting too far.)}
\end{truyen}
\begin{baihoc}
Dừng lại để nhìn mình không phải chậm trễ; đó là một phần của trưởng thành.
\textit{(Pausing to examine yourself is not delay; it is part of maturity.)}
\end{baihoc}
\begin{ghinhoanh}
\textbf{Short English to remember:}

Slow steps can see more.
\end{ghinhoanh}
\begin{tuhocnhanh}
1. Lâu rồi em đã dừng lại để nhìn mình chưa? \textit{(How long has it been since you paused to look at yourself?)}

2. Nếu chậm lại hôm nay, em sẽ thấy gì rõ hơn? \textit{(If you slowed down today, what would become clearer?)}
\end{tuhocnhanh}
\ngancach
